\chapter{总结与展望}\label{chap:conclusion}

\section{工作总结}
本文围绕轻量级实时数字协作白板系统 QuickBoard 的设计与实现开展研究与工程实践，面向短会话协作场景，强调无需注册、快速创建/加入房间的使用体验。系统以三服务架构实现前端白板、后端协作与公式识别能力：前端基于 Vue~3 与 Fabric.js 提供对象级绘制与编辑；后端基于 Node.js（Express）与 Socket.IO 提供房间管理、实时事件分发与服务端权威裁决；OCR 服务基于 FastAPI 提供公式识别能力，并通过后端代理统一错误码与 traceId。

围绕本课题的核心工程风险点，本文完成并验证了以下工作：
\begin{itemize}
  \item \textbf{一致性策略落地}：采用服务端权威裁决与字段级 LWW 合并，结合 \texttt{roomVersion} 增量补包、快照兜底与手动对齐，形成可解释、可对齐的协作路径。评估表明该策略在丢包与断线窗口条件下能够实现最终收敛。
  \item \textbf{性能与降载优化}：在渲染侧通过视口裁剪与对象缓存降低大对象规模下的平均渲染耗时；在网络侧通过 Ghost Brush 将高频过程预览与权威状态解耦，并通过批处理与限包显著降低消息数量。
  \item \textbf{公式识别闭环}：实现“框选—裁剪—识别—确认插入—可编辑渲染”的闭环交互，并通过统一错误码、traceId 与可选 debug 回传提升失败可定位性，使公式能力在教学场景下具备可用性与可维护性。
\end{itemize}

综合测试与评估结果，QuickBoard 在工程复杂度可控的前提下实现了轻量协作白板的核心能力，并通过对照实验与基准测试提供了可复核的论证证据。

\section{后续展望}
后续工作可从以下方向进一步完善系统：
\begin{itemize}
  \item \textbf{一致性语义增强}：在保持系统可解释性的前提下，针对极端并发同字段编辑的意图保持进行增强；可探索在局部数据结构上引入 CRDT/OT 思路，或对高冲突对象设计更细粒度的数据模型。
  \item \textbf{持久化与历史回放}：在轻量前提下完善房间历史与版本回溯，使教学场景能够沉淀推导过程并支持复习与复盘。
  \item \textbf{权限与审计}：增加房间访问控制、主持人权限与操作审计，提升公开演示与课堂使用中的可控性。
  \item \textbf{OCR 量化评估与迭代}：构建更系统的样本集与量化指标（成功率、耗时分位数、可编辑率），并在不牺牲隐私与延迟的前提下持续提升识别质量。
  \item \textbf{跨端体验与可访问性}：完善快捷键、触控适配与视觉反馈，并针对移动端与不同分辨率设备持续优化交互与性能表现。
\end{itemize}
